\vspace*{4cm}
\begin{center}
    \begin{Arabic}
        {\Huge\bfseries الملخص}
    \end{Arabic}
\end{center}
\vskip 1cm

\phantomsection
\addcontentsline{toc}{chapter}{\texorpdfstring{\textarabic{الملخص}}{الملخص}}

\begin{Arabic}
أدت التحولات الرقمية في قطاع الصحة إلى ظهور الأجهزة الطبية المتصلة، والسجلات الصحية الإلكترونية،
ومنصات المراقبة عن بعد، وأنظمة دعم القرار المعتمدة على الذكاء الاصطناعي. ورغم أن هذه التقنيات
تنتج كميات كبيرة من البيانات المتنوعة، فإن العديد من أنظمة الصحة الإلكترونية ما تزال مجزأة، وثابتة،
ومحدودة التخصيص. وتقدم تقنية التوأم الرقمي مقاربة واعدة لمعالجة هذه الحدود، من خلال إنشاء
تمثيلات افتراضية ديناميكية للكيانات الفيزيائية أو العمليات أو المرضى، يتم تحديثها باستمرار عبر تبادل
البيانات، وتدعيمها بالمحاكاة والتحليل وخدمات دعم القرار.

يعرض هذا البحث مراجعة منظمة لحالة الفن حول توظيف التوائم الرقمية في مجال الصحة الإلكترونية. ويتناول
أولا أسس التوأم الرقمي، بما في ذلك التمييز بين النموذج الرقمي، والظل الرقمي، والتوأم الرقمي،
والتوأم الرقمي الذكي، ثم يقدم بيئة الصحة الإلكترونية والصحة الرقمية التي يمكن تطبيق هذه المبادئ
داخلها. بعد ذلك، يحلل تطبيقات التوائم الرقمية الموجهة للرعاية الصحية، ومعمارياتها، والتقنيات الممكنة
لها، وتحديات تنفيذها. وتستعمل التوائم الرقمية القلبية الوعائية كحالة سريرية ممثلة لأنها تجمع بين
الإشارات الحيوية، والتصوير الطبي، والقياسات الهيموديناميكية، والمتابعة الطولية، والنمذجة الخاصة
بالمريض. كما يولي البحث
اهتماما خاصا لقضايا قابلية التشغيل البيني، وحوكمة البيانات، والخصوصية، والأمن، وقابلية التفسير،
والتحقق، والاعتبارات الأخلاقية. وتستخلص مراجعة الأدبيات معمارية عامة متعددة الطبقات للتوائم الرقمية
في الصحة الإلكترونية، إضافة إلى مبادئ تتعلق بتدفق البيانات، والمتطلبات العابرة للطبقات، والتقييم.
أما الفصل الأخير فيلخص دراسة الماستر ويقدم مشروع CardioTwin كاستمرار عملي في المجال القلبي
الوعائي ضمن مشروع نهاية الدراسة، مع الحفاظ على الطابع العام للبحث حول التوائم الرقمية للصحة
الإلكترونية.

\textbf{الكلمات المفتاحية:} التوأم الرقمي، الصحة الإلكترونية، الصحة الرقمية، الرعاية الصحية،
إنترنت الأشياء، إنترنت الأشياء الطبية، الذكاء الاصطناعي، المحاكاة، قابلية التشغيل البيني، دعم القرار
السريري، التوأم الرقمي القلبي الوعائي.
\end{Arabic}

\clearpage
